\section{Kitty Problem y principio abierto/cerrado}

\subsection{Enunciado}
En su libro \emph{Fundamentals of Software Architecture}, Richards y Ford presentan el «Kitty
Problem» para ilustrar un fallo común en el diseño de arquitecturas orientadas a servicios.
Describa este problema: un requisito transversal (como la entrega de gatitos) afecta a múltiples
servicios previamente existentes (TaxiUI, TaxiFinder, TaxiSelector, etc.), requiriendo cambios
coordinados en todos ellos.

¿Por qué la solución de dividir el sistema en servicios basados en verbos o entidades
funcionales (como TaxiFinder o TaxiSelector) llevó a este problema? Basándose en los principios
de DDD y la Arquitectura Onion/Clean, ¿cómo se podría haber diseñado el sistema para que un
cambio de este tipo afectara a un solo componente o servicio? Proponga una solución utilizando
componentes como Kittens y Rides que se añadan como plugins, explicando cómo esto respeta el
Principio de Abierto/Cerrado (OCP) y qué papel juegan los Adaptadores en esta extensión.

\subsection{Respuesta}
\textit{[Redactar con \parencite{richards2020fundamentals,martin2017clean,palermo2008onion}.]}
